\documentclass[10pt,a4paper]{amsart}
\usepackage[margin=19mm]{geometry}
\usepackage{amsmath,amssymb,mathtools}
\usepackage[scr=rsfs,cal=euler]{mathalfa}
\usepackage{xfrac}
\usepackage[unicode,hidelinks,bookmarks=false]{hyperref}
\newcommand{\ie}{i.e.\ }
\newcommand{\cf}{cf.\ }
\newcommand{\resp}{respectively\ }
\newcommand{\Subsection}{\subsection}
\providecommand{\nobreakdash}{\nobreak\hbox{-}}
\setlength{\emergencystretch}{2em}
\multlinegap0pt
\usepackage{amssymb}
\newtheorem{theorem}{Theorem}
\title{Geometry and Topology of Gradient Shrinking Sasaki-Ricci Solitons}
\author{Shu-Cheng Chang, Yingbo Han, Chin-Tung Wu}
\date{Source: arXiv:2508.13495v2 (CC BY 4.0)}
\begin{document}
\maketitle

\noindent\emph{Source and licence.} Geometry and Topology of Gradient Shrinking Sasaki-Ricci Solitons. Version arXiv:2508.13495v2 (CC BY 4.0), distributed by arXiv under the Creative Commons Attribution 4.0 International licence, http://creativecommons.org/licenses/by/4.0/, as recorded in the arXiv licence metadata for this exact version and shown on the abstract page https://arxiv.org/abs/2508.13495v2. Full licence notice: \texttt{licenses/arxiv-2508.13495-CC-BY-4.0.txt}.\par\smallskip

\noindent\textit{Editorial scope.} Counted unit: the article's theorem Thm2 (source0.tex lines 216--237). The article's complete original proof of it is delivered with it (source0.tex lines 239--517, one \texttt{proof} environment of 279 lines). No mathematics is written, repaired or re-derived: the statement and the proof are the article's own lines, in the article's own notation, and no retained line is substituted or re-flowed.  No further source-local context is needed: every cross-reference of every retained line resolves inside the two delivered spans.  Nothing was omitted from any retained span, so the package carries no omission record.  No retained line contains a citation, so the wrapper carries no bibliography and no bibliography entry.  No retained line contains a figure environment, an image-inclusion command or drawing code, so no image is delivered and the package declares no asset dependency.  Editorial formatting only, in addition: an AMS \texttt{amsart} preamble that restates the article's own declarations and macros used by the retained lines, namely the environments \texttt{theorem} on separate counters, together with the standard packages those lines need.  The retained environments print in this document as Theorem 1 in order of appearance, whereas the article prints the same items with the numbering of its own full document.  The complete native source of the article as distributed in the arXiv e-print for this exact version is bundled unchanged as \texttt{sources/183481/source0.tex}.
\begin{theorem}
\label{Thm2}For a Sasakian manifold $(M,\xi,\eta,g,\Phi)$. We assume that
$\varphi:M\rightarrow\mathbb{R}$ is a basic function which satisfies
$\varphi_{\alpha\beta}=0$. Additionally, there exists a constant $C>0$ such
that
\[%
\begin{array}
[c]{c}%
|\nabla\varphi|\leq C(d(x_{0},x)+1)
\end{array}
\]
for any $x\in M$. Then any real solution $u$ to
\[%
\begin{array}
[c]{c}%
\Delta_{\varphi}u:=\Delta u-\left\langle \nabla u,\nabla\varphi\right\rangle
=0\text{\ \textrm{with}\ }\int_{M}|\nabla u|^{2}e^{-\varphi}<\infty
\end{array}
\]
must be transversely plurihamonic and $u_{00}=0$. Moreover, it is constant if
$\varphi$ is proper on $M$.
\end{theorem}

\begin{proof}
For a fixed point $x_{0}\in M$, let the cut-off function $\phi$ be defined by%
\[%
\begin{array}
[c]{c}%
\phi(x)=\left\{
\begin{array}
[c]{ll}%
1 & \mathrm{on}\text{ }B_{x_{0}}(R)\\
\frac{2R-d(x_{0},x)}{R} & \mathrm{on}\text{ }B_{x_{0}}(2R)\backslash B_{x_{0}%
}(R)\\
0 & \mathrm{on}\text{ }M\backslash B_{x_{0}}(2R).
\end{array}
\right.
\end{array}
\]
We then consider the following integration for second derivatives of $u$
\begin{equation}%
\begin{array}
[c]{ll}
& \int_{M}\left[  |u_{\alpha\bar{\beta}}|^{2}+|u_{0\alpha}|^{2}+|u_{0\bar
{\alpha}}|^{2}+|u_{00}|^{2}\right]  \phi^{2}e^{-\varphi}\\
= & \int_{M}u_{\alpha\bar{\beta}}u_{\bar{\alpha}\beta}\phi^{2}e^{-\varphi
}+u_{0\alpha}u_{0\bar{\alpha}}\phi^{2}e^{-\varphi}+u_{0\bar{\alpha}}%
u_{0\alpha}\phi^{2}e^{-\varphi}+u_{00}u_{00}\phi^{2}e^{-\varphi},
\end{array}
\label{11}%
\end{equation}
here $u_{0}=\xi u.$ We investigate each term in (\ref{11}). Integration by
parts implies
\begin{equation}%
\begin{array}
[c]{ll}
& \int_{M}u_{\alpha\bar{\beta}}u_{\bar{\alpha}\beta}\phi^{2}e^{-\varphi}\\
= & \int_{M}\mathrm{\operatorname{Re}}(u_{\alpha\bar{\beta}}[u_{\beta
\bar{\alpha}}+\sqrt{-1}g_{\alpha\bar{\beta}}u_{0}])\phi^{2}e^{-\varphi}%
=\int_{M}\mathrm{\operatorname{Re}}(u_{\alpha\bar{\beta}}u_{\beta\bar{\alpha}%
})\phi^{2}e^{-\varphi}\\
= & -\int_{M}\mathrm{\operatorname{Re}}(u_{\alpha\bar{\beta}\bar{\alpha}%
}u_{\beta})\phi^{2}e^{-\varphi}-\int_{M}\mathrm{\operatorname{Re}}%
(u_{\alpha\bar{\beta}}u_{\beta}[(\phi^{2})_{\bar{\alpha}}-\varphi_{\bar
{\alpha}}\phi^{2}])e^{-\varphi}\\
= & -\int_{M}\mathrm{\operatorname{Re}}(u_{\alpha\bar{\alpha}\bar{\beta}%
}u_{\beta})\phi^{2}e^{-\varphi}-\int_{M}\mathrm{\operatorname{Re}}%
(u_{\alpha\bar{\beta}}u_{\beta}[(\phi^{2})_{\bar{\alpha}}-\varphi_{\bar
{\alpha}}\phi^{2}])e^{-\varphi}\\
= & \int_{M}\mathrm{\operatorname{Re}}([u_{00}-\Delta u]_{\bar{\beta}}%
u_{\beta})\phi^{2}e^{-\varphi}-\int_{M}\mathrm{\operatorname{Re}}%
(u_{\alpha\bar{\beta}}u_{\beta}[(\phi^{2})_{\bar{\alpha}}-\varphi_{\bar
{\alpha}}\phi^{2}])e^{-\varphi}\\
= & \int_{M}\mathrm{\operatorname{Re}}(u_{00\bar{\beta}}u_{\beta}-(\Delta
u)_{\bar{\beta}}u_{\beta})\phi^{2}e^{-\varphi}-\int_{M}%
\mathrm{\operatorname{Re}}(u_{\alpha\bar{\beta}}u_{\beta}[(\phi^{2}%
)_{\bar{\alpha}}-\varphi_{\bar{\alpha}}\phi^{2}])e^{-\varphi}.
\end{array}
\label{12}%
\end{equation}
here $\Delta u=u_{\alpha\bar{\alpha}}+u_{00}$ and $\mathrm{\operatorname{Re}%
}(z):=\frac{z+\overline{z}}{2}$ denotes the real part of the complex number
$z.$ Integration by parts again
\begin{equation}%
\begin{array}
[c]{ll}
& \int_{M}u_{0\alpha}u_{0\bar{\alpha}}\phi^{2}e^{-\varphi}=\int_{M}%
\mathrm{\operatorname{Re}}(u_{\alpha0}u_{0\bar{\alpha}})\phi^{2}e^{-\varphi}\\
= & -\int_{M}\mathrm{\operatorname{Re}}(u_{0\bar{\alpha}0}u_{\alpha})\phi
^{2}e^{-\varphi}-\int_{M}\mathrm{\operatorname{Re}}(u_{\bar{\alpha}0}%
u_{\alpha})(\phi^{2})_{0}e^{-\varphi}\\
= & -\int_{M}\mathrm{\operatorname{Re}}([u_{00\bar{\alpha}}+u_{\bar{\alpha}%
}]u_{\alpha})\phi^{2}e^{-\varphi}-\int_{M}\mathrm{\operatorname{Re}}%
(u_{\bar{\alpha}0}u_{\alpha})(\phi^{2})_{0}e^{-\varphi},
\end{array}
\label{13}%
\end{equation}
here we used $\varphi$ is basic, that is $\varphi_{0}=0.$ Also
\begin{equation}%
\begin{array}
[c]{ll}
& \int_{M}u_{0\bar{\alpha}}u_{\alpha0}\phi^{2}e^{-\varphi}\\
= & -\int_{M}\mathrm{\operatorname{Re}}(u_{\alpha0\bar{\alpha}})u_{0}\phi
^{2}e^{-\varphi}-\int_{M}\mathrm{\operatorname{Re}}(u_{\alpha0}[(\phi
^{2})_{\bar{\alpha}}-\varphi_{\bar{\alpha}}\phi^{2}])u_{0}e^{-\varphi}\\
= & -\int_{M}u_{\alpha\bar{\alpha}0}u_{0}\phi^{2}e^{-\varphi}-\int
_{M}\mathrm{\operatorname{Re}}(u_{\alpha0}[(\phi^{2})_{\bar{\alpha}}%
-\varphi_{\bar{\alpha}}\phi^{2}])u_{0}e^{-\varphi}.
\end{array}
\label{14}%
\end{equation}
The last term in (\ref{11}) becomes%
\begin{equation}%
\begin{array}
[c]{c}%
\int_{M}u_{00}u_{00}\phi^{2}e^{-\varphi}=-\int_{M}u_{000}u_{0}\phi
^{2}e^{-\varphi}-\int_{M}u_{00}u_{0}(\phi^{2})_{0}e^{-\varphi}.
\end{array}
\label{15}%
\end{equation}
By combining (\ref{12}), (\ref{13}), (\ref{14}) and (\ref{15}) with
(\ref{11}), we obtain%
\begin{equation}%
\begin{array}
[c]{ll}
& \int_{M}\left[  |u_{\alpha\bar{\beta}}|^{2}+|u_{0\alpha}|^{2}+|u_{0\bar
{\alpha}}|^{2}+|u_{00}|^{2}\right]  \phi^{2}e^{-\varphi}\\
\leq & -\int_{M}\left[  \mathrm{\operatorname{Re}}((\Delta u)_{\bar{\beta}%
}u_{\beta}+(\Delta u)_{0}u_{0})\right]  \phi^{2}e^{-\varphi}+\int
_{M}\operatorname{Re}(u_{\alpha\bar{\beta}}u_{\beta}\varphi_{\bar{\alpha}%
})\phi^{2}e^{-\varphi}\\
& +\int_{M}\operatorname{Re}(u_{\alpha0}u_{0}\varphi_{\bar{\alpha}})\phi
^{2}e^{-\varphi}-\int_{M}[\mathrm{\operatorname{Re}}(u_{\bar{\alpha}%
0}u_{\alpha})+u_{00}u_{0}](\phi^{2})_{0}e^{-\varphi}\\
& -\int_{M}\operatorname{Re}((u_{\alpha\bar{\beta}}u_{\beta}+u_{\alpha0}%
u_{0})(\phi^{2})_{\bar{\alpha}})e^{-\varphi}.
\end{array}
\label{16}%
\end{equation}


First, we deal with the first term in (\ref{16}),
\[%
\begin{array}
[c]{ll}
& -\int_{M}\mathrm{\operatorname{Re}}((\Delta u)_{\bar{\beta}}u_{\beta
}+(\Delta u)_{0}u_{0})\phi^{2}e^{-\varphi}=-\int_{M}\left\langle \nabla\Delta
u,\nabla u\right\rangle \phi^{2}e^{-\varphi}\\
= & \int_{M}(\Delta u)(\Delta_{\varphi}u)\phi^{2}e^{-\varphi}+\int_{M}(\Delta
u)\left\langle \nabla u,\nabla\phi^{2}\right\rangle e^{-\varphi}\\
= & \int_{M}(\Delta u)\left\langle \nabla u,\nabla\phi^{2}\right\rangle
e^{-\varphi},
\end{array}
\]
thus by the Cauchy-Schwarz inequality
\begin{equation}%
\begin{array}
[c]{ll}
& -\int_{M}\mathrm{\operatorname{Re}}((\Delta u)_{\bar{\beta}}u_{\beta
}+(\Delta u)_{0}u_{0})\phi^{2}e^{-\varphi}\\
\leq & \frac{1}{8(2n+1)}\int_{M}(\Delta u)^{2}\phi^{2}e^{-\varphi}%
+8(2n+1)\int_{M}|\nabla u|^{2}|\nabla\phi|^{2}e^{-\varphi}\\
\leq & \frac{1}{4}\int_{M}[|u_{\alpha\bar{\beta}}|^{2}+|u_{00}|^{2}]\phi
^{2}e^{-\varphi}+8(2n+1)\int_{M}|\nabla u|^{2}|\nabla\phi|^{2}e^{-\varphi}.
\end{array}
\label{17}%
\end{equation}
Secondly integration by parts yields
\[%
\begin{array}
[c]{ll}
& \int_{M}\operatorname{Re}(u_{\alpha\bar{\beta}}u_{\beta}\varphi_{\bar
{\alpha}})\phi^{2}e^{-\varphi}\\
= & -\int_{M}\operatorname{Re}(u_{\alpha}u_{\beta\overline{\beta}}%
\varphi_{\bar{\alpha}})\phi^{2}e^{-\varphi}-\int_{M}\operatorname{Re}%
(u_{\alpha}u_{\beta}\varphi_{\bar{\alpha}\overline{\beta}})\phi^{2}%
e^{-\varphi}\\
& +\int_{M}\operatorname{Re}(u_{\alpha}u_{\beta}\varphi_{\bar{\alpha}}%
\varphi_{\overline{\beta}})\phi^{2}e^{-\varphi}-\int_{M}\operatorname{Re}%
(u_{\alpha}u_{\beta}\varphi_{\bar{\alpha}}(\phi^{2})_{\overline{\beta}%
})e^{-\varphi},
\end{array}
\]
where we have used the assumption $\varphi_{\alpha\beta}=\varphi_{\bar{\alpha
}\overline{\beta}}=0.$ Thus%
\begin{equation}%
\begin{array}
[c]{ll}
& \int_{M}\operatorname{Re}(u_{\alpha\bar{\beta}}u_{\beta}\varphi_{\bar
{\alpha}})\phi^{2}e^{-\varphi}\\
= & -\int_{M}(\Delta u)\operatorname{Re}(u_{\alpha}\varphi_{\bar{\alpha}}%
)\phi^{2}e^{-\varphi}+\int_{M}\operatorname{Re}(u_{\alpha}\varphi_{\bar
{\alpha}})u_{00}\phi^{2}e^{-\varphi}\\
& +\int_{M}\operatorname{Re}(u_{\alpha}u_{\beta}\varphi_{\bar{\alpha}}%
\varphi_{\overline{\beta}})\phi^{2}e^{-\varphi}-\int_{M}\operatorname{Re}%
(u_{\alpha}u_{\beta}\varphi_{\bar{\alpha}}(\phi^{2})_{\overline{\beta}%
})e^{-\varphi}\\
= & -\int_{M}\langle\nabla u,\nabla\varphi\rangle^{2}\phi^{2}e^{-\varphi}%
+\int_{M}\operatorname{Re}(u_{\alpha}\varphi_{\bar{\alpha}})u_{00}\phi
^{2}e^{-\varphi}\\
& +\int_{M}\operatorname{Re}(u_{\alpha}u_{\beta}\varphi_{\bar{\alpha}}%
\varphi_{\overline{\beta}})\phi^{2}e^{-\varphi}-\int_{M}\operatorname{Re}%
(u_{\alpha}u_{\beta}\varphi_{\bar{\alpha}}(\phi^{2})_{\overline{\beta}%
})e^{-\varphi}.
\end{array}
\label{18}%
\end{equation}
Note that%
\[%
\begin{array}
[c]{c}%
\operatorname{Re}(u_{\alpha}u_{\beta}\varphi_{\bar{\alpha}}\varphi
_{\overline{\beta}})\leq\langle\nabla u,\nabla\varphi\rangle^{2}.
\end{array}
\]
From the first term in line 3 in (\ref{16}), we have
\[%
\begin{array}
[c]{c}%
\int_{M}\operatorname{Re}(u_{\alpha0}u_{0}\varphi_{\bar{\alpha}})\phi
^{2}e^{-\varphi}=-\int_{M}\operatorname{Re}(u_{\alpha}\varphi_{\bar{\alpha}%
})u_{00}\phi^{2}e^{-\varphi}-\int_{M}\mathrm{\operatorname{Re}}(u_{\alpha
}\varphi_{\bar{\alpha}})u_{0}(\phi^{2})_{0}e^{-\varphi},
\end{array}
\]
since $\varphi$ is basic, so $\varphi_{\bar{\alpha}0}=\varphi_{0\bar{\alpha}%
}=0.$ Plugging these into (\ref{18}), we conclude%
\begin{equation}%
\begin{array}
[c]{c}%
\int_{M}\operatorname{Re}(u_{\alpha\bar{\beta}}u_{\beta}\varphi_{\bar{\alpha}%
}+u_{\alpha0}u_{0}\varphi_{\bar{\alpha}})\phi^{2}e^{-\varphi}\leq4\int
_{M}|\nabla u|^{2}|\nabla\varphi|\phi|\nabla\phi|e^{-\varphi}.
\end{array}
\label{19}%
\end{equation}
For the last two terms in (\ref{16}). By the Cauchy-Schwarz inequality again%
\begin{equation}%
\begin{array}
[c]{ll}
& -\int_{M}[\mathrm{\operatorname{Re}}(u_{\bar{\alpha}0}u_{\alpha}%
)+u_{00}u_{0}](\phi^{2})_{0}+\operatorname{Re}((u_{\alpha\bar{\beta}}u_{\beta
}+u_{\alpha0}u_{0})(\phi^{2})_{\bar{\alpha}})]e^{-\varphi}\\
\leq & \frac{1}{4}\int_{M}[|u_{\alpha\bar{\beta}}|^{2}+2|u_{\alpha0}%
|^{2}+2|u_{\bar{\alpha}0}|^{2}+|u_{00}|]\phi^{2}e^{-\varphi}\\
& +c_{1}(n)\int_{M}|\nabla u|^{2}|\nabla\phi|^{2}e^{-\varphi}+c_{2}(n)\int
_{M}|\nabla\varphi||\nabla u|^{2}\phi|\nabla\phi|e^{-\varphi}.
\end{array}
\label{19a}%
\end{equation}
Therefore, by substituting (\ref{17}), (\ref{19}) and (\ref{19a}) into
(\ref{16}), we final get
\begin{equation}%
\begin{array}
[c]{ll}
& \int_{M}\left[  |u_{\alpha\bar{\beta}}|^{2}+|u_{0\alpha}|^{2}+|u_{0\bar
{\alpha}}|^{2}+|u_{00}|^{2}\right]  \phi^{2}e^{-\varphi}\\
\leq & c_{3}(n)\int_{M}|\nabla u|^{2}|\nabla\phi|^{2}e^{-\varphi}+c_{4}%
(n)\int_{M}|\nabla\varphi||\nabla u|^{2}\phi|\nabla\phi|e^{-\varphi}.
\end{array}
\label{19b}%
\end{equation}
Since $\int_{M}|\nabla u|^{2}e^{-\varphi}<\infty$, it implies that%
\[%
\begin{array}
[c]{c}%
\int_{M}|\nabla u|^{2}|\nabla\phi|^{2}e^{-\varphi}\rightarrow0\text{
}\mathrm{as}\text{ }R\rightarrow\infty.
\end{array}
\]
Furthermore, by the assumption that $|\nabla\varphi|\leq C(d(x_{0},x)+1)$, we
have $|\nabla\varphi||\nabla\phi|\leq C$ on $M$. It yields that
\[%
\begin{array}
[c]{c}%
\int_{M}|\nabla\varphi||\nabla u|^{2}\phi|\nabla\phi|e^{-\varphi}%
\rightarrow0\text{ }\mathrm{as}\text{ }R\rightarrow\infty.
\end{array}
\]
We conclude that $u_{\alpha\bar{\beta}}=0,$ $u_{\alpha0}=0$ and $u_{00}=0$ by
letting $R\rightarrow\infty$ in (\ref{19b}).

To show $u$ is constant when $\varphi$ is proper. Note that $\langle\nabla
u,\nabla\varphi\rangle=0$ as $u$ is both $\varphi$-harmonic and harmonic. Let
\[%
\begin{array}
[c]{c}%
D(t)=\{x:\varphi(x)\leq t\},
\end{array}
\]
whish is compact as $\varphi$ is proper. Now we compute
\[%
\begin{array}
[c]{c}%
\int_{D(t)}|\nabla u|^{2}=\frac{1}{2}\int_{D(t)}\Delta u^{2}=\frac{1}{2}%
\int_{\partial D(t)}\frac{\partial u^{2}}{\partial\nu}=\int_{\partial
D(t)}u\frac{\langle\nabla u,\nabla\varphi\rangle}{|\nabla\varphi|}=0.
\end{array}
\]
Since this is true for any $t$, it follows that $|\nabla u|=0$ on $M$, so $u$
is constant.
\end{proof}

\end{document}
